\documentclass[11pt,letterpaper]{article}
\usepackage[margin=0.7in,headheight=15pt]{geometry}
\usepackage[T1]{fontenc}\usepackage{lmodern,amsmath,amssymb,fancyhdr,parskip,hyperref}
\hypersetup{colorlinks=true,urlcolor=blue}
\pdfinfoomitdate=1\pdftrailerid{}\pdfsuppressptexinfo=15
\pagestyle{fancy}\fancyhf{}\lhead{Week 83 / Hackenbush balances}\rhead{Facilitator draft}\cfoot{Bellingham Math Circle / Week 83 / facilitator / \thepage}
\setlength{\parskip}{7pt}\setlength{\parindent}{0pt}
\begin{document}
\section*{Mathematical overview}
Under normal play Blue cuts one blue segment and Red cuts one red segment; everything disconnected from the ground disappears. Only one component changes per move, and a player with no move loses. Finite red/blue stalks have dyadic number values. Read stalk strings \emph{bottom to top}: B is 1, R is $-1$, BR is $1/2$, BRR is $1/4$, and RB is $-1/2$.

Two BR stalks plus a separate R have value 0, as do two BRR stalks plus a separate RB. In a zero game the second player can win whichever color starts; there is no draw. Positive/negative value means Blue/Red has a winning strategy for either starter. Numeric equality is equivalence in \emph{all} disjunctive-sum contexts, not equality of edge counts or observed duration. Finite play on selected boards motivates the known equality theorem; it does not prove every context.

Optional Grades 4--5 nursery: choose the earliest-born number strictly above all left bounds and below all right bounds. This needs every left bound below every right bound; then exactly one earliest-born number fits. So $\{0\mid3\}=1$, $\{-1\mid1\}=0$, $\{0\mid1\}=1/2$, $\{1/2\mid1\}=3/4$. Finite birthdays produce dyadics. The rule is not arbitrary midpoint insertion. The same rule gives every stalk its number: put the numbers of the boards Blue's cuts leave on the left and Red's on the right. So $\mathrm{BR}=\{0\mid1\}=1/2$, $\mathrm{BRR}=\{0\mid 1/2,1\}=1/4$ and $\mathrm{RB}=\{-1\mid0\}=-1/2$ (Problem 6), which is why the balanced boards of Problems 2 and 3 sum to 0.

What children are meant to discover: a stalk's worth is not its number of segments. Two BR balance one R (Problem 2); two BRR balance one RB, but two BR do not (Problem 3); a stalk beside its colour swap balances (Problem 4). Playing a board with each starter is the experiment, ``this board balances'' is a conjecture, and a winning reply to every opening is the explanation. Band map: Grades 2--3 play Problems 1--4 and find the balances; Grades 4--5 also explain every reply and, on a return visit, compute stalk numbers in the nursery (Problems 5--6). K--1 has no pages; an adult-led cube route is on page 2.
\newpage
\section*{Preparation and launch}
Per pair: 12 blue and 12 red snap cubes (interlocking cubes that join into towers) and one student packet. Build each printed stalk as a tower standing on the table, bottom segment first; the table is the ground. A cut pulls off the chosen cube together with every cube above it, and all of them leave the board, so the cubes themselves carry out ``all segments cut off from the ground disappear.'' The largest boards need 3 blue and 5 red cubes (Problem 3C) and 3 blue and 4 red (Problem 2C). If a child has trouble telling the colours apart, mark the cubes B and R with stickers. Every stalk here is a simple rooted path; no floating component counts as a move. Fallback without cubes: sleeved boards and erasable pens, with the partner checking that the fallen part is crossed out. Cube towers are unrehearsed with this packet.

Prerequisites: alternate turns, distinguish colors, preserve support after a cut. Fractions are unnecessary to start. Optional nursery additionally needs signed numbers, halves/quarters and strict inequalities, plus reading birthday cards. There are no K--1 pages; the K--1 route below uses cubes and an adult.

After handling materials, launch with the non-target worked cut on page 1: build B,R,B as a tower, let Red pull off the R, and show that the top B comes away with it. The partner checks that only grounded segments remain. Replay once with the other color. Show how separate stalks share one turn rule before distributing balances.

Cut the fifteen birthday-number cards on page 7 only for the optional nursery; use page 8 as the reusable invention board. First route: 10 minutes easy starter-swapped games; 15 minutes balance candidates; 10 minutes building new balanced boards (Problem 4); five minutes sharing response strategies. The nursery is a return visit, or a choice for children already fluent. Adult at each table; upper triplet rotates Blue/Red/referee. Use the same workspace rather than page-by-page copying.

K--1 route (adult-led, no K--1 pages). The adult reads the rules and builds Problem 1's boards as cube towers; two children each take a colour and play every board twice, swapping who starts. Then ask: ``Make a board where whoever goes second wins,'' using towers of one colour only. Key in plain words: with one-colour towers a player can always pull off just the top cube, and taking more only wastes turns, so the colour with more cubes wins; with as many blue cubes as red (B + R, BB + RR, B + B + RR) whoever goes second wins. That K--1 children can keep track of turns and starters is a prediction, not an observation.
\newpage
\section*{Exact game reasoning and solution methods}
The cut rules give $\mathrm{BR}=\{0\mid1\}=1/2$ and $\mathrm{BRR}=\{0\mid1,1/2\}=1/4$; removing the upper red leaves BR, removing the lower red leaves B. The number assigned is the simplest between the effective bounds. Reversing colors negates the value. Disjoint grounded stalks add as games.

Winners and winning moves, one line per board (values in parentheses):
\begin{itemize}
\item 1A, B (1): Blue; with Red starting, Red has no move. 1B, R ($-1$): Red. 1C, B + R (0): second player; no colour always wins.
\item 1D, BB + R (1): Blue, but Blue starting must cut the top B; cutting the bottom B takes both blue segments and Red wins. Red starting cuts R, then Blue cuts either B.
\item 2A, BR + R ($-1/2$): Red. Red starting cuts the top R of BR, leaving B + R with Blue to move. Blue starting can only remove BR; Red cuts R.
\item 2B, BR + BR + R (0): second player; the proof follows.
\item 2C, BR + BR + BR + R ($1/2$): Blue. Blue starting removes a whole BR, leaving 2B with Red to move. If Red starts on the lone R, Blue removes a whole BR; if Red cuts a top R, Blue removes the B left behind or a whole BR.
\item 3A, BRR + RB ($-1/4$): Red. Red starting must cut the top R of BRR, leaving BR + RB (0). If Blue removes BRR, Red cuts the R of RB; if Blue cuts the top B of RB, Red cuts either R of BRR.
\item 3B, BR + BR + RB ($1/2$): Blue, though it looks like 2B. Blue starting wins with any cut: removing a BR leaves BR + RB, cutting RB's top B leaves BR + BR + R, both 0. If Red cuts a top R, Blue removes the B left behind; if Red cuts the R of RB, Blue removes a whole BR.
\item 3C, BRR + BRR + RB (0): second player. If Blue removes a BRR, Red cuts the other's top R (BR + RB); if Blue cuts RB's top B, Red cuts a top R ($-1/4$). If Red cuts a top or middle R, Blue removes the other whole BRR; if Red cuts the R of RB, Blue removes a whole BRR.
\end{itemize}
The balance $\mathrm{BR}+\mathrm{BR}+\mathrm{R}$ has value 0. Here is a complete concrete response proof:
\begin{itemize}
\item Blue starts by cutting a BR base. Red cuts the other BR's top R, leaving B+R with Blue to move. Blue removes B, and Red removes R last.
\item Red starts by cutting the lone R. Blue removes one whole BR at its base; Red must cut the remaining top R, then Blue removes its base last.
\item Red starts by cutting a BR top R. Blue removes the other whole BR, leaving B+R with Red to move. Red removes R, then Blue removes B last.
\end{itemize}
Symmetric choices of the two identical BR stalks need no new case. A strategy covers every opening; one successful play is evidence only.

In Problem 3, $\mathrm{BRR}+\mathrm{BRR}+\mathrm{RB}=1/4+1/4-1/2=0$ balances, while $\mathrm{BR}+\mathrm{BR}+\mathrm{RB}=1/2$ does not: RB is worth half of R, so the pattern ``two stalks balance one'' from Problem 2 fails. The balance is board C, not B, so the position on the page does not answer the question. Verify printed boards against the independent stalk solver included in the source package.

If students count edges as scores, ask them to play B versus BR and compare a cut at the bottom. If they only try one starter, invite the other starter before inferring a winner class. Do not call a second-player win a tie. Green edges, if later introduced, give the first-player game star rather than the number zero; they are not part of the core packet.
\newpage
\section*{Nursery: simplicity, not a midpoint procedure}
Day 0 gives 0. Day 1 introduces $-1,1$. Day 2 introduces $-2,-1/2,1/2,2$. Day 3 introduces\[ -3,-3/2,-3/4,-1/4,1/4,3/4,3/2,3. \] Cards record \emph{first} birthdays; an older number does not acquire a new birthday when represented by new constraints.

Problem 4 asks for boards not in Problems 1--3. Examples: BR + RB, and any stalk beside its colour swap, since the second player can copy each cut on the other stalk; B + RB + RB; BB + RR. Problem 5 gives $1,0,1/2,3/4$ in printed order.

Problem 6: on BR, Blue's cut leaves the empty board (0) and Red's leaves B (1), so BR is the earliest card between 0 and 1, $1/2$ (day 2). On BRR, Blue's cut leaves 0, and Red's cuts leave B (1) or BR ($1/2$), so BRR is the earliest card above 0 and below $1/2$, that is $1/4$ (day 3). On RB, Blue's cut leaves R ($-1$) and Red's leaves the empty board (0), so RB is $-1/2$. Children need BR's card before BRR's. With these, Problem 2B is $1/2+1/2-1=0$ and Problem 3C is $1/4+1/4-1/2=0$.

A left-bound card means strictly greater; right means strictly smaller. Select the earliest birthday among permitted cards. $\{0\mid3\}$ admits the day-1 number 1, so the midpoint $3/2$ is not the answer. $\{-1\mid1\}$ describes the already-existing number 0. $\{1/2\mid1\}$ first admits $3/4$ on day 3. When the left card is below the right card but no card fits (for example $1/4$ and $1/2$, whose number $3/8$ is born on day 4), say ``not in this catalog'' before discussing later days. When the left card is not below the right card, no number of any day fits.

Optional adult doorway: no finite-day card is greater than all $0,1,2,\ldots$; the simplest surreal with those lower bounds is $\omega$. A number strictly above 0 and below $1,1/2,1/4,\ldots$ leads to a positive surreal infinitesimal. These are mathematical descriptions of infinite bound sets, not physically performed days.

Keep operations distinct: ordinary queue concatenation in Week 82 has $1+\omega=\omega$, while surreal field addition is commutative and gives $1+\omega=\omega+1$. Naming both operations ``plus'' without their meaning would teach a false transfer. No transfinite symbol manipulation is required on student pages.
\newpage
\section*{Sources, encore and observation record}
Source: Dierk Schleicher and Michael Stoll, \emph{An Introduction to Conway's Games and Numbers}, arXiv math/0410026v2, Sections 2--3, especially the number construction and finite-birthday dyadic characterization. \url{https://arxiv.org/html/math/0410026v2}. The game semantics and arithmetic are established mathematics; our examples, cut illustrations and printed boards are original teaching designs. No external worksheet is copied.

Nursery return visit: find a pair of bounds, one card on each side, that chooses each card; every card except $-3$ and $3$ can be chosen this way. Or make two different pairs of bounds that choose the same card, and a pair that chooses a card newer than both bounds.

Encore: build two different stalk collections with the same value and challenge a partner to distinguish them by adding a small test stalk. Explain why successful small tests suggest, but do not establish, full game equality. A later visit can introduce a single green edge and investigate why its behavior cannot be any number.

Pedagogy: Rozhkovskaya, Lesson 7 ``At the lesson,'' item 1, observed missed legal moves in a game. The referee and cut/drop demonstration are our response; Lessons 3/8 support visible possibilities and lighter transcription. This is not classroom validation of Hackenbush.

Unpiloted. US Letter, single-sided, Actual Size. Cube towers, color legibility on a particular printer, preparation and classroom procedure untested. Record which board/starter combinations were actually played, whether children removed fallen pieces, and whether balance explanations covered all replies. Preserve observations separately from hypotheses about abstraction readiness.
\end{document}
